% =====================================================================
%  Conclusion : importée par main.tex avec % =====================================================================
%  Conclusion : importée par main.tex avec % =====================================================================
%  Conclusion : importée par main.tex avec % =====================================================================
%  Conclusion : importée par main.tex avec \include{conclusion}.
%  Section non numérotée, ajoutée au sommaire.
% =====================================================================

% Macros de Dirac (déjà définies dans main.tex ; ce bloc évite l'erreur
% « Undefined control sequence \ket » si main.tex est une ancienne version).
\providecommand{\ket}[1]{|#1\rangle}
\providecommand{\bra}[1]{\langle #1|}
\providecommand{\braket}[2]{\langle #1|#2\rangle}

\phantomsection
\section*{Conclusion}
\addcontentsline{toc}{section}{Conclusion}

On est parti de l'onde de matière de de Broglie et de l'équation de Schrödinger pour
aboutir à un modèle de calcul complet. Le tableau ci-dessous résume la traduction de la
physique ondulatoire vers l'informatique quantique.

\begin{center}
\renewcommand{\arraystretch}{1.45}
\begin{tabular}{@{}p{3.6cm}p{5.4cm}p{5.6cm}@{}}
\toprule
\textbf{Notion} & \textbf{Mécanique ondulatoire} & \textbf{Informatique quantique} \\
\midrule
État & fonction d'onde $\Psi(x,t)$ & vecteur $\ket{\Psi}$ d'un espace de Hilbert, ou $\rho$ \\
Grandeur physique & opérateur $\hat p$, $\hat H$ & matrice hermitienne \\
Évolution & équation de Schrödinger & porte unitaire $U$ \\
Mesure & règle de Born $|\Psi|^2$ & $P(i)=|\braket{i}{\Psi}|^2$ et projection \\
Plusieurs systèmes & fonction de plusieurs variables & produit tensoriel $\otimes$ \\
Corrélations & --- & intrication, états de Bell \\
\bottomrule
\end{tabular}
\end{center}

L'intrication est la ressource qui distingue les systèmes quantiques : le protocole de
téléportation du chapitre~\ref{chap:bell} la consomme pour transmettre un état inconnu, sans le copier.

\paragraph{Pour aller plus loin.} Ce document ne traite ni les algorithmes (Deutsch, Grover,
Shor \parencite{shor1997}), ni les inégalités de Bell et CHSH \parencite{bell1964}, ni la
correction d'erreurs quantiques, qui prolongent directement les chapitres~\ref{chap:portes}
et~\ref{chap:bell}. Les ouvrages de \textcite{nielsen2010} et de \textcite{mermin2007} en
donnent une présentation complète.
.
%  Section non numérotée, ajoutée au sommaire.
% =====================================================================

% Macros de Dirac (déjà définies dans main.tex ; ce bloc évite l'erreur
% « Undefined control sequence \ket » si main.tex est une ancienne version).
\providecommand{\ket}[1]{|#1\rangle}
\providecommand{\bra}[1]{\langle #1|}
\providecommand{\braket}[2]{\langle #1|#2\rangle}

\phantomsection
\section*{Conclusion}
\addcontentsline{toc}{section}{Conclusion}

On est parti de l'onde de matière de de Broglie et de l'équation de Schrödinger pour
aboutir à un modèle de calcul complet. Le tableau ci-dessous résume la traduction de la
physique ondulatoire vers l'informatique quantique.

\begin{center}
\renewcommand{\arraystretch}{1.45}
\begin{tabular}{@{}p{3.6cm}p{5.4cm}p{5.6cm}@{}}
\toprule
\textbf{Notion} & \textbf{Mécanique ondulatoire} & \textbf{Informatique quantique} \\
\midrule
État & fonction d'onde $\Psi(x,t)$ & vecteur $\ket{\Psi}$ d'un espace de Hilbert, ou $\rho$ \\
Grandeur physique & opérateur $\hat p$, $\hat H$ & matrice hermitienne \\
Évolution & équation de Schrödinger & porte unitaire $U$ \\
Mesure & règle de Born $|\Psi|^2$ & $P(i)=|\braket{i}{\Psi}|^2$ et projection \\
Plusieurs systèmes & fonction de plusieurs variables & produit tensoriel $\otimes$ \\
Corrélations & --- & intrication, états de Bell \\
\bottomrule
\end{tabular}
\end{center}

L'intrication est la ressource qui distingue les systèmes quantiques : le protocole de
téléportation du chapitre~\ref{chap:bell} la consomme pour transmettre un état inconnu, sans le copier.

\paragraph{Pour aller plus loin.} Ce document ne traite ni les algorithmes (Deutsch, Grover,
Shor \parencite{shor1997}), ni les inégalités de Bell et CHSH \parencite{bell1964}, ni la
correction d'erreurs quantiques, qui prolongent directement les chapitres~\ref{chap:portes}
et~\ref{chap:bell}. Les ouvrages de \textcite{nielsen2010} et de \textcite{mermin2007} en
donnent une présentation complète.
.
%  Section non numérotée, ajoutée au sommaire.
% =====================================================================

% Macros de Dirac (déjà définies dans main.tex ; ce bloc évite l'erreur
% « Undefined control sequence \ket » si main.tex est une ancienne version).
\providecommand{\ket}[1]{|#1\rangle}
\providecommand{\bra}[1]{\langle #1|}
\providecommand{\braket}[2]{\langle #1|#2\rangle}

\phantomsection
\section*{Conclusion}
\addcontentsline{toc}{section}{Conclusion}

On est parti de l'onde de matière de de Broglie et de l'équation de Schrödinger pour
aboutir à un modèle de calcul complet. Le tableau ci-dessous résume la traduction de la
physique ondulatoire vers l'informatique quantique.

\begin{center}
\renewcommand{\arraystretch}{1.45}
\begin{tabular}{@{}p{3.6cm}p{5.4cm}p{5.6cm}@{}}
\toprule
\textbf{Notion} & \textbf{Mécanique ondulatoire} & \textbf{Informatique quantique} \\
\midrule
État & fonction d'onde $\Psi(x,t)$ & vecteur $\ket{\Psi}$ d'un espace de Hilbert, ou $\rho$ \\
Grandeur physique & opérateur $\hat p$, $\hat H$ & matrice hermitienne \\
Évolution & équation de Schrödinger & porte unitaire $U$ \\
Mesure & règle de Born $|\Psi|^2$ & $P(i)=|\braket{i}{\Psi}|^2$ et projection \\
Plusieurs systèmes & fonction de plusieurs variables & produit tensoriel $\otimes$ \\
Corrélations & --- & intrication, états de Bell \\
\bottomrule
\end{tabular}
\end{center}

L'intrication est la ressource qui distingue les systèmes quantiques : le protocole de
téléportation du chapitre~\ref{chap:bell} la consomme pour transmettre un état inconnu, sans le copier.

\paragraph{Pour aller plus loin.} Ce document ne traite ni les algorithmes (Deutsch, Grover,
Shor \parencite{shor1997}), ni les inégalités de Bell et CHSH \parencite{bell1964}, ni la
correction d'erreurs quantiques, qui prolongent directement les chapitres~\ref{chap:portes}
et~\ref{chap:bell}. Les ouvrages de \textcite{nielsen2010} et de \textcite{mermin2007} en
donnent une présentation complète.
.
%  Section non numérotée, ajoutée au sommaire.
% =====================================================================

% Macros de Dirac (déjà définies dans main.tex ; ce bloc évite l'erreur
% « Undefined control sequence \ket » si main.tex est une ancienne version).
\providecommand{\ket}[1]{|#1\rangle}
\providecommand{\bra}[1]{\langle #1|}
\providecommand{\braket}[2]{\langle #1|#2\rangle}

\phantomsection
\section*{Conclusion}
\addcontentsline{toc}{section}{Conclusion}

On est parti de l'onde de matière de de Broglie et de l'équation de Schrödinger pour
aboutir à un modèle de calcul complet. Le tableau ci-dessous résume la traduction de la
physique ondulatoire vers l'informatique quantique.

\begin{center}
\renewcommand{\arraystretch}{1.45}
\begin{tabular}{@{}p{3.6cm}p{5.4cm}p{5.6cm}@{}}
\toprule
\textbf{Notion} & \textbf{Mécanique ondulatoire} & \textbf{Informatique quantique} \\
\midrule
État & fonction d'onde $\Psi(x,t)$ & vecteur $\ket{\Psi}$ d'un espace de Hilbert, ou $\rho$ \\
Grandeur physique & opérateur $\hat p$, $\hat H$ & matrice hermitienne \\
Évolution & équation de Schrödinger & porte unitaire $U$ \\
Mesure & règle de Born $|\Psi|^2$ & $P(i)=|\braket{i}{\Psi}|^2$ et projection \\
Plusieurs systèmes & fonction de plusieurs variables & produit tensoriel $\otimes$ \\
Corrélations & --- & intrication, états de Bell \\
\bottomrule
\end{tabular}
\end{center}

L'intrication est la ressource qui distingue les systèmes quantiques : le protocole de
téléportation du chapitre~\ref{chap:bell} la consomme pour transmettre un état inconnu, sans le copier.

\paragraph{Pour aller plus loin.} Ce document ne traite ni les algorithmes (Deutsch, Grover,
Shor \parencite{shor1997}), ni les inégalités de Bell et CHSH \parencite{bell1964}, ni la
correction d'erreurs quantiques, qui prolongent directement les chapitres~\ref{chap:portes}
et~\ref{chap:bell}. Les ouvrages de \textcite{nielsen2010} et de \textcite{mermin2007} en
donnent une présentation complète.
